\documentclass{article}
\usepackage[utf8]{inputenc}
\usepackage{geometry}
\usepackage{booktabs}
\usepackage{amsmath}
\usepackage{pifont}
\usepackage{enumitem}
\usepackage{tabularx}

\geometry{a4paper, margin=1in}

\newcommand{\cmark}{\ding{51}}%
\newcommand{\xmark}{\ding{55}}%

\title{How ORPO Works: Preference Alignment and Quality Gating}
\date{}

\begin{document}

\maketitle

\section{Introduction and Definitions}

\subsection{WHAT IS ORPO?}
\begin{itemize}
    \item \textbf{ORPO = Odds Ratio Preference Optimization}. It is a preference-optimization method for fine-tuning language models.
    \item ORPO combines the supervised fine-tuning (SFT) objective with an odds-ratio preference objective in a single training stage.
    \item It learns from preferred and rejected responses without requiring a separately trained reward model.
\end{itemize}

\subsection{WHAT IS A PREFERENCE PAIR?}
A preference pair contains a prompt and two responses to that prompt: one preferred and one rejected. These pairs are prepared as training data before ORPO fine-tuning.
\begin{itemize}
    \item[\cmark] \textbf{Preferred answer} (higher quality / better aligned)
    \item[\xmark] \textbf{Rejected answer} (lower quality / less suitable)
\end{itemize}

\noindent\textit{Illustrative example:} Prompt: ``Explain photosynthesis in simple words.'' Preferred response: ``Plants use sunlight to make food from water and carbon dioxide.'' Rejected response: ``Photosynthesis is the process by which atoms form nuclear explosions.'' These are illustrative preference-training examples, not guaranteed outputs produced by ORPO.

\raggedbottom
\newpage

\section{End-to-End Pipeline: From Preference Pairs to\\ Runtime Quality Gating}

\subsection{Training Stage}

\begin{enumerate}
    \item \textbf{Prepare preference data}
    \begin{itemize}
        \item Each record contains a prompt, a preferred response, and a rejected response.
        \item The prompt and its two candidate responses form a preference pair for training.
    \end{itemize}
    
    \item \textbf{Provide the preference pair}
    \begin{itemize}
        \item \textit{Prompt:} ``Explain photosynthesis in simple words.''
        \item[\cmark] \textbf{Preferred:} Plants use sunlight to make food from water and carbon dioxide.
        \item[\xmark] \textbf{Rejected:} Photosynthesis is the process by which atoms form nuclear explosions.
    \end{itemize}
    
    \item \textbf{Apply single-stage ORPO fine-tuning}
    \begin{itemize}
        \item ORPO combines the supervised learning loss with an odds-ratio preference loss for the pair $(y_w, y_l)$, where $y_w$ is preferred and $y_l$ is rejected.
        \item In simplified form, its objective can be written as $\mathcal{L}_{\mathrm{ORPO}} = \mathcal{L}_{\mathrm{SFT}} + \lambda\mathcal{L}_{\mathrm{OR}}$. It does not use a separate reward signal or reward model.
    \end{itemize}
    
    \item \textbf{Produce the fine-tuned model}
    \begin{itemize}
        \item The result is a fine-tuned model that has learned from the supplied preference data.
    \end{itemize}
\end{enumerate}

\subsection{Runtime Inference Stage}

\begin{enumerate}
    \item \textbf{Receive a new user prompt} and generate a response using the fine-tuned model.
    \item \textbf{Apply the separate quality gate} to the generated response.
    \begin{itemize}
        \item \textbf{Example checks:} task relevance, coherence, and applicable safety requirements.
        \item The quality gate is an optional project-level runtime component; it is not an inherent part of ORPO.
    \end{itemize}
    
    \item \textbf{Approve or trigger an action}
    \begin{itemize}
        \item[\cmark] \textbf{Approved:} The response passes the defined checks.
        \item[\xmark] \textbf{Rejected or revised:} The system blocks the response or triggers a revision according to the defined rules.
    \end{itemize}
\end{enumerate}

\newpage

\section{Comparison and Selection Evaluation}

\subsection{ORPO vs. Traditional PPO-Based RLHF}

\begin{table}[h!]
\centering
\caption{Comparison of ORPO and Traditional PPO-Based RLHF}
\begin{tabularx}{\textwidth}{@{}X X@{}}
\toprule
\textbf{ORPO} & \textbf{Traditional PPO-Based RLHF} \\ \midrule
\cmark\ Combines SFT and preference optimization in one training stage. & \cmark\ Typically uses SFT, a separately trained reward model, and PPO optimization. \\
\cmark\ Does not require a separate reward model. & \cmark\ Requires a separate reward model to provide the learned objective for PPO. \\
\cmark\ Has fewer separate training components. & \cmark\ Has a more involved multi-stage training pipeline. \\
\bottomrule
\end{tabularx}
\end{table}

\subsection{Why Choose ORPO for This Project?}
\begin{itemize}
    \item \textbf{Direct preference learning:} it learns from preferred and rejected responses without a separate reward-model training step.
    \item \textbf{Simpler training workflow:} combining supervised fine-tuning and preference optimization reduces the number of separate training components to implement and explain.
    \item \textbf{Additional runtime validation:} the separate quality gate can check relevance, coherence, and applicable safety requirements after generation. ORPO itself does not guarantee that every response is safe or correct.
\end{itemize}

\vfill
\footnotesize
\noindent\rule{\textwidth}{0.4pt}
\begin{center}
Group 01 \textbullet\ CSE, CEK \hfill Fine-tuning, ORPO \hfill ORPO
\end{center}

\end{document}
